\section{Related Work}
\label{sec:related_work}

Our work is situated at the intersection of model merging, test-time adaptation, and parameter-efficient channel routing. We provide a comprehensive overview of each of these areas, emphasizing how our proposed EdgeMerge addresses the core practical limitations of existing research.

\subsection{Parameter-Space Model Merging}
Model merging has become a key technique in modern deep learning pipelines, enabling developers to combine multiple fine-tuned models into a single multi-task model without retraining. The simplest approach, linear weight averaging, was popularized by Model Soups \cite{wortsman2022model} and Task Arithmetic \cite{ilharco2023editing}. Task Arithmetic shows that adding fine-tuning weight updates (task vectors) to a pre-trained base model creates a unified model capable of executing multiple downstream tasks.

To combat the performance degradation caused by inter-task weight conflicts, several static weight-space consolidation methods have been proposed. Fisher Weighted Averaging \cite{matena2022merging} uses diagonal Fisher Information matrices to weigh parameter averages based on their task-specific sensitivities. TIES-Merging \cite{yadav2023resolving} resolves parameter interference by pruning redundant task-vector elements, performing a sign-consensus check to resolve directional conflicts, and averaging only the agreed-upon updates. Other approaches like Git Re-Basin \cite{ainsworth2023git} and ZipIt! \cite{stoica2023zipit} align the internal representations of disparate models by resolving permutation symmetries before averaging weights. While these methods are computationally cheap and preserve the original model size, they are static and uniform, ignoring the fact that task-specific weights should be modulated dynamically depending on their local visual or textual context.

\subsection{Adaptive Merging and Test-Time Adaptation}
To introduce flexibility, recent papers have proposed optimizing merging coefficients dynamically. AdaMerging \cite{yu2024adamerging} optimizes layer-wise merging coefficients by backpropagating losses computed on unlabeled test batches. SyMerge \cite{lu2026symerge} introduces an agentic low-rank adapter scaling framework that adjusts layer-wise scales based on the outputs of a pool of task-expert teacher networks at test-time. FoldMerge maps model weights into a latent "Origami Space" using normalizing flows and optimizes the coordinate system over 500 gradient steps.

These approaches are closely related to Test-Time Adaptation (TTA) frameworks, such as Tent \cite{wang2021tent} and Memo \cite{zhang2022memo}, which update model parameters or normalization layers on-the-fly using entropy minimization or self-consistency objectives on test streams.

From an engineering and deployment perspective, these adaptive merging and TTA techniques suffer from severe computational overhead. Running gradient descent at test-time introduces substantial backpropagation latency (often requiring minutes of optimization on high-end GPUs) and extremely high peak memory consumption due to the need to store activation graphs and multiple teacher networks in memory. Furthermore, as noted in recent literature, optimizing delicate parameters on a tiny, transductive test-time batch is highly prone to overfitting, which degrades the model's performance on future unseen data. EdgeMerge resolves these critical issues by completely eliminating the backpropagation process, extracting adaptive channel-wise coefficients in a single, near-instant forward pass.

\subsection{Activation-Based Saliency and Channel Modulation}
Our method is conceptually inspired by activation-based neural network compression and modulation. In channel pruning \cite{he2017channel, li2016pruning, molchanov2016pruning}, the importance of individual convolutional or linear filters is estimated using activation scales or first-order Taylor expansions on a calibration set, allowing developers to discard redundant channels. In dynamic network routing, Squeeze-and-Excitation networks \cite{hu2018squeeze} and Convolutional Block Attention Modules (CBAM) \cite{woo2018cbam} introduce dynamic channel gating at inference time by passing activations through lightweight bottleneck layers.

While traditional channel gating requires modifying the neural network architecture and training new parameters from scratch, EdgeMerge applies activation-based channel routing directly to model weights post-hoc. By analyzing the scale-normalized delta activations between each task expert and the base model, we can compute channel-wise salience in closed-form. This allows us to perform fine-grained channel gating on the visual projection layer directly within the merged parameter space, introducing zero architectural modifications or inference latency.
